\chapter{Il front-end: HTML}

\section{Il browser come ambiente di esecuzione}

Il client-side di una web app \`e costituito da: una o pi\`u \textbf{pagine HTML}, ciascuna
corredata da uno o pi\`u \textbf{fogli di stile CSS} e associata a zero o pi\`u
\textbf{moduli JavaScript} (una pagina senza JS \`e passiva; in un'app multi-pagina alcune
pagine possono essere di sola lettura).

\begin{importante}[Sandbox]
Il browser vede \emph{ogni pagina come a s\'e stante}: ciascuna viene eseguita in una
``scatola chiusa'' (\textbf{sandbox}) creata al caricamento. Se la pagina viene ricaricata,
la sandbox viene distrutta e ricreata: \textbf{la memoria dell'applicazione si perde}.
Si pu\`o ovviare parzialmente con i meccanismi di \emph{storage locale} del browser,
condivisibili fra pagine della stessa \emph{origine} (stesso server), ma con diversi limiti.
Se il client-side ha pi\`u pagine, \`e compito del \emph{server-side} gestire il flusso di
informazioni fra una pagina e l'altra.
\end{importante}

Ricevuti HTML + CSS + JS, il browser usa HTML e CSS per inizializzare la \emph{vista}:
\begin{itemize}
  \item il \textbf{documento HTML} definisce la \emph{struttura logica} e la \emph{semantica}
        delle parti della pagina (sezioni, menu di navigazione, ruolo dei testi, media,
        elementi di input, collegamenti ipertestuali);
  \item i \textbf{fogli CSS} sono insiemi di regole che definiscono l'\emph{aspetto grafico}
        degli elementi (caratteri, colori, dimensioni, collocazione);
  \item con HTML e CSS il browser inizializza il \textbf{DOM} (Document Object Model): la
        struttura dati che mantiene in memoria tutte le informazioni sulla visualizzazione;
        pu\`o essere modificata durante la vita della pagina, che termina alla chiusura o al
        ricaricamento.
\end{itemize}

\section{HTML: un linguaggio di annotazione}

HTML descrive struttura e aspetto di un documento multimediale e ipertestuale usando solo
testo semplice (codifica UTF-8) e le URL dei media da includere. Ogni \textbf{elemento} \`e
caratterizzato da:
\begin{itemize}
  \item un \textbf{tipo};
  \item un \textbf{contenuto} (testo e/o altri elementi; alcuni tipi non lo prevedono);
  \item una serie di \textbf{attributi} opzionali, dipendenti dal tipo.
\end{itemize}
La \emph{specifica formale} di HTML definisce i tipi esistenti, il loro ``significato
inteso'', il contenuto atteso e gli attributi possibili. Il browser usa la strutturazione
(ad albero) per organizzare la visualizzazione.

\begin{esempio}[Sintassi di un elemento]
\begin{lstlisting}[language=HTMLCSS]
<XYZ attrib=val> contenuto </XYZ>
\end{lstlisting}
Tag di \emph{apertura} (nome del tipo fra parentesi angolari), \emph{contenuto} (testo e/o
altri elementi), tag di \emph{chiusura} (uguale ma con \texttt{/}). Gli attributi si indicano
nel tag di apertura. Attributi non previsti dalla specifica sono ignorati dal browser, ma
funzioni JavaScript possono sfruttare attributi \emph{custom} definiti dal programmatore,
che per convenzione hanno nome della forma \texttt{data-*}.
\end{esempio}

\subsection{Struttura di un documento}
\begin{lstlisting}[language=HTMLCSS]
<!DOCTYPE html>
<html lang="en">
  <head> ... </head>
  <body> ... </body>
</html>
\end{lstlisting}
\begin{itemize}
  \item \textbf{Preambolo} \texttt{<!DOCTYPE html>}: dichiara la versione della specifica;
        esiste per retrocompatibilit\`a ed \`e di fatto sempre uguale.
  \item \textbf{HEAD} (intestazione): \emph{meta-informazioni} sul documento, usate dal
        browser ma non visualizzate come contenuti.
  \item \textbf{BODY} (corpo): il contenuto effettivo, ci\`o che viene visualizzato.
\end{itemize}

\subsection{Elementi a blocco e in linea}
In prima approssimazione gli elementi si dividono in \textbf{elementi a blocco} (definiscono
un riquadro) ed \textbf{elementi in linea} (seguono il fluire del testo). In un documento ben
strutturato gli elementi a blocco stanno solo dentro altri elementi a blocco: un blocco dentro
un elemento in linea porta a comportamenti imprevisti. Il browser ``fa del suo meglio'' per
visualizzare comunque, ma specie progettando UI serve una visualizzazione prevedibile e
controllabile, non affidata alle strategie compensative del browser.

\section{Collegamenti ipertestuali}
\begin{lstlisting}[language=HTMLCSS]
<a href="https://it.wikipedia.org/wiki/...">testo cliccabile</a>
\end{lstlisting}
\begin{itemize}
  \item L'elemento \`e \textbf{A} (\emph{anchor}); pu\`o contenere sia elementi in linea che
        a blocco. Il contenuto \`e lo \textbf{hotspot}: ci\`o su cui l'utente clicca.
  \item La destinazione (\emph{target}) \`e nell'attributo \textbf{href}, una URL valida:
        \textbf{esterna} (specifica dominio e path, dominio diverso dall'attuale),
        \textbf{assoluta} (stesso dominio, path completo dalla radice, inizia con \texttt{/}),
        \textbf{relativa} (stesso dominio, path relativo al documento attuale).
\end{itemize}

\section{Gli elementi principali}

\subsection{Testo}
Elementi \emph{a blocco}:
\begin{itemize}
  \item \textbf{P}: paragrafo.
  \item \textbf{H1\dots H6}: titoli a importanza decrescente.
  \item \textbf{UL}: lista a punti; contiene elementi \textbf{LI}. Attributo \texttt{type}
        (\texttt{circle}, \texttt{disc}, \texttt{square}) per il tipo di punto.
  \item \textbf{OL}: lista numerata di \textbf{LI}. \texttt{type} (\texttt{a}, \texttt{A},
        \texttt{i}, \texttt{I}, \texttt{1}) per numerazione alfabetica/romana/intera;
        \texttt{start} per il valore iniziale.
\end{itemize}
Elementi \emph{in linea}:
\begin{itemize}
  \item \textbf{EM}: testo da enfatizzare (default corsivo).
  \item \textbf{STRONG}: enfasi forte (default grassetto).
  \item \textbf{SPAN}: testo da distinguere; il browser di default non lo distingue,
        l'aspetto si stabilisce poi via CSS.
\end{itemize}

\subsection{Tabelle}
Per contenuti realmente tabellari (dati in righe/colonne).
\begin{attenzione}
\textbf{Non usare le tabelle per il layout della pagina}: per la disposizione dei contenuti
esistono gli elementi strutturali in combinazione con i CSS.
\end{attenzione}
\begin{itemize}
  \item \textbf{TABLE}: l'intera tabella; contiene righe.
  \item \textbf{TR}: una riga; contiene celle.
  \item \textbf{TD}: una cella. Attributi \texttt{rowspan}/\texttt{colspan} (interi $>0$,
        default 1) per celle che spaziano su pi\`u righe/colonne; la somma dei colspan
        dovrebbe essere costante per riga.
  \item \textbf{TH}: cella di intestazione; \texttt{scope} = \texttt{row} o \texttt{col}.
  \item \textbf{THEAD, TBODY, TFOOT}: raggruppano le righe in intestazione, corpo, pi\`e.
\end{itemize}

\subsection{Contenuti multimediali}
\begin{itemize}
  \item \textbf{IMG}: immagine, elemento \emph{in linea}. \texttt{src} = URL (esterna,
        assoluta o relativa); \texttt{width}/\texttt{height} consigliati, cos\`i il browser
        imposta la pagina senza attendere il download.
  \item \textbf{AUDIO}: player audio incastonato (\emph{embedded}); \texttt{src},
        \texttt{autoplay} (true/false), \texttt{controls} (mostrare i comandi).
  \item \textbf{VIDEO}: analogo, con anche \texttt{width}/\texttt{height}.
\end{itemize}
\begin{attenzione}[Ridimensionare non riduce il traffico]
L'immagine viene comunque scaricata \emph{nella dimensione originale} e solo dopo
ridimensionata dal browser. Se l'obiettivo \`e limitare il traffico di rete, servono pi\`u
versioni dell'immagine sul server, usando la URL di quella della dimensione voluta.
\end{attenzione}

\subsection{Elementi strutturali}
L'uso varia fra pagine ``redazionali'' (documenti, articoli, blog: interazione minima) e
``applicative'' (la pagina \`e la UI di un'applicazione: utente attivo).
\begin{itemize}
  \item \textbf{DIV}: suddivisione strutturale \emph{senza significato} a priori;
        fondamentale per dare un layout preciso, quasi indispensabile nelle UI.
  \item \textbf{SECTION}: suddivide la pagina in sezioni indipendenti.
  \item \textbf{ARTICLE}: contraddistingue un articolo presentabile anche indipendentemente
        dal resto.
  \item \textbf{NAV}: area dedicata a un menu di navigazione.
  \item \textbf{HEADER}: intestazione \emph{visibile} di pagina/sezione (titolo, logo).
        \textbf{Da non confondere con HEAD}, che contiene la meta-intestazione non
        visualizzata.
  \item \textbf{FOOTER}: pi\`e di pagina (contatti, cookie, ecc.).
\end{itemize}

\section{I moduli (form) HTML}

L'HTML puro permette moduli con cui l'utente invia dati al web server (sensato solo se il
server li riceve con una app in grado di elaborarli). Qui parliamo di form \emph{gestiti dal
browser}; vedremo poi form gestiti da JavaScript, pi\`u flessibili.

\subsection{Il tag FORM}
Definisce il comportamento del modulo tramite gli attributi:
\begin{itemize}
  \item \textbf{action}: la URL (esterna, relativa o assoluta) a cui inviare i dati;
  \item \textbf{method}: il method HTTP. Con \textbf{GET} i dati vanno codificati come
        \emph{query string} nella URL; con \textbf{POST} vanno nel \emph{body} con
        \texttt{Content-Type: application/x-www-form-urlencoded}.
\end{itemize}
In entrambi i casi i dati hanno la forma
\texttt{nome1=valore1\&nome2=valore2\&\dots\&nomeK=valoreK}: i nomi vengono dagli elementi
di input, i valori da ci\`o che l'utente ha inserito.

\subsection{Gli elementi di input}
Dentro un FORM pu\`o esserci qualsiasi HTML; per raccogliere dati servono elementi di input,
per lo pi\`u col tag \textbf{INPUT} e il suo attributo \texttt{type}.
\begin{attenzione}
Tutti gli INPUT richiedono \emph{obbligatoriamente} l'attributo \textbf{name}: \`e il nome
del campo codificato nella richiesta HTTP.
\end{attenzione}
\begin{itemize}
  \item \texttt{<input type="text"/>}: casella di testo. Attributi: \texttt{size}
        (larghezza in caratteri), \texttt{value} (valore preimpostato modificabile),
        \texttt{placeholder} (indicazione che \emph{non} conta come valore e scompare
        digitando).
  \item \texttt{<input type="checkbox"/>}: \texttt{value} = valore riportato se
        contrassegnata; \texttt{checked} = inizialmente selezionata.
  \item \texttt{<input type="radio"/>}: i radio-button di uno stesso gruppo devono avere lo
        \emph{stesso} \texttt{name} (solo uno selezionabile); \texttt{value} identifica quale
        \`e stato scelto; \texttt{checked} per la selezione iniziale.
  \item \texttt{<input type="submit" value="Premi qui"/>}: \textbf{obbligatorio} in un form
        HTML puro; non fornisce un dato, ma provoca l'invio secondo \texttt{action} e
        \texttt{method}; \texttt{value} \`e il testo sul pulsante.
  \item \texttt{<input type="reset" value="Ricomincia"/>}: opzionale; riporta tutti gli input
        allo stato iniziale.
\end{itemize}
Altri due elementi comuni nei form:
\begin{itemize}
  \item \textbf{LABEL}: etichetta esplicativa associata a un INPUT (essenziale per checkbox e
        radio). L'INPUT deve avere un \texttt{id}; lo stesso valore va nell'attributo
        \texttt{for} della LABEL: \texttt{<label for="inputid">\dots</label>}.
  \item \textbf{SELECT}: menu a tendina; \texttt{name} sta sul SELECT. Le opzioni sono
        elementi \textbf{OPTION} al suo interno, ciascuna con il proprio \texttt{value}
        (riportato nella richiesta in base alla scelta):
\begin{lstlisting}[language=HTMLCSS]
<select name="scelta">
  <option value="opt1">Prima opzione</option>
  <option value="opt2">Seconda opzione</option>
</select>
\end{lstlisting}
\end{itemize}
